\documentclass[9pt,a4paper,twocolumn]{ctexart}

% --- Packages ---
\usepackage[top=1.2cm, bottom=1.2cm, left=1.2cm, right=1.2cm, columnsep=0.8cm]{geometry}
\usepackage{hyperref}
\usepackage{graphicx}
\usepackage{float}
\usepackage{longtable}
\usepackage{tabularx}
\usepackage{booktabs}
\usepackage{enumitem}
\usepackage{xcolor}
\usepackage{fancyhdr}
\usepackage{listings}
\usepackage{fontspec}
\usepackage{titlesec}
\usepackage{caption}

% --- Code block styling ---
\definecolor{codebg}{RGB}{245,245,245}
\definecolor{codeframe}{RGB}{220,220,220}
\definecolor{codekw}{RGB}{0,0,192}
\definecolor{codecomment}{RGB}{0,128,0}
\definecolor{codestring}{RGB}{163,21,21}
\definecolor{codenum}{RGB}{128,0,128}
\definecolor{codepunct}{RGB}{90,90,90}
\definecolor{badgegreen}{RGB}{34,139,34}
\definecolor{badgeblue}{RGB}{30,144,255}
\definecolor{badgeblack}{RGB}{50,50,50}

% --- Difficulty badges (rounded corners via tikz, pre-rendered in saveboxes) ---
\usepackage{tikz}
\newsavebox{\badgechubox}
\savebox{\badgechubox}{\tikz[baseline]{\node[fill=badgegreen,text=white,rounded corners=2.5pt,inner xsep=4pt,inner ysep=1.5pt,font=\sffamily\footnotesize\bfseries]{初级};}}
\newsavebox{\badgezhongbox}
\savebox{\badgezhongbox}{\tikz[baseline]{\node[fill=badgeblue,text=white,rounded corners=2.5pt,inner xsep=4pt,inner ysep=1.5pt,font=\sffamily\footnotesize\bfseries]{中级};}}
\newsavebox{\badgegaobox}
\savebox{\badgegaobox}{\tikz[baseline]{\node[fill=badgeblack,text=white,rounded corners=2.5pt,inner xsep=4pt,inner ysep=1.5pt,font=\sffamily\footnotesize\bfseries]{高级};}}
\newcommand{\lvchu}{\usebox{\badgechubox}}
\newcommand{\lvzhong}{\usebox{\badgezhongbox}}
\newcommand{\lvgao}{\usebox{\badgegaobox}}
\pdfstringdefDisableCommands{%
  \def\lvchu{[初级]}%
  \def\lvzhong{[中级]}%
  \def\lvgao{[高级]}%
}

\lstset{
  basicstyle=\ttfamily\scriptsize,
  keywordstyle=\color{codekw}\bfseries,
  commentstyle=\color{codecomment}\itshape,
  stringstyle=\color{codestring},
  backgroundcolor=\color{codebg},
  frame=single,
  rulecolor=\color{codeframe},
  breaklines=true,
  breakatwhitespace=false,
  breakindent=0pt,
  postbreak=\mbox{\textcolor{red}{$\hookrightarrow$}\space},
  numbers=left,
  numberstyle=\tiny\color{gray},
  columns=flexible,
  keepspaces=true,
  showstringspaces=false,
  tabsize=2,
  aboveskip=4pt,
  belowskip=4pt,
  extendedchars=true,
  literate={，}{{，}}1 {；}{{；}}1 {！}{{！}}1 {？}{{？}}1 {“}{{“}}1 {”}{{”}}1 {（}{{（}}1 {）}{{）}}1 {《}{{《}}1 {》}{{》}}1,
}

% --- Extra language definitions (no built-in listings support) ---
\lstdefinelanguage{json}{
  morekeywords={true,false,null},
  sensitive=true,
  morestring=[b]",
  comment=[l]{//},
  morecomment=[s]{/*}{*/},
}
\lstdefinelanguage{yaml}{
  morekeywords={true,false,null,yes,no,on,off},
  sensitive=false,
  morecomment=[l]{\#},
  morestring=[b]",
  morestring=[d]',
}
\lstdefinelanguage{http}{
  morekeywords={GET,POST,PUT,DELETE,PATCH,HEAD,OPTIONS,HTTP,Host,Accept,Authorization,Content-Type,Content-Length,User-Agent},
  sensitive=true,
  morecomment=[l]{\#},
  morestring=[b]",
}
\lstdefinelanguage{TypeScript}{
  morekeywords={abstract,any,as,async,await,break,case,catch,class,const,continue,debugger,declare,default,delete,do,else,enum,export,extends,false,finally,for,from,function,get,if,implements,import,in,instanceof,interface,keyof,let,module,namespace,never,new,null,number,of,package,private,protected,public,readonly,return,set,static,string,super,switch,symbol,this,throw,true,try,type,typeof,undefined,unknown,var,void,while,with,yield,Record,Array,Promise,AsyncIterable,ReadableStream,Date,Error,Object,JSON},
  sensitive=true,
  morecomment=[l]{//},
  morecomment=[s]{/*}{*/},
  morestring=[b]",
  morestring=[b]',
  morestring=[b]`{`},
}

% --- Compact spacing ---
\linespread{0.95}
\setlength{\parskip}{1pt plus 1pt minus 1pt}
\setlength{\parindent}{0pt}

% --- Table styling ---
\renewcommand{\arraystretch}{1.1}

% --- Heading styling (numbered) ---
\titleformat{\section}{\normalsize\bfseries}{\thesection\quad}{0em}{}
\titlespacing{\section}{0pt}{6pt}{2pt}
\titleformat{\subsection}{\small\bfseries}{\thesubsection\quad}{0em}{}
\titlespacing{\subsection}{0pt}{4pt}{1pt}
\titleformat{\subsubsection}{\small\bfseries}{\thesubsubsection\quad}{0em}{}
\titlespacing{\subsubsection}{0pt}{3pt}{1pt}

% --- Page style ---
\pagestyle{fancy}
\fancyhf{}
\fancyhead[L]{\small\emph{\leftmark}}
\fancyhead[R]{\small\thepage}
\renewcommand{\headrulewidth}{0.4pt}
\renewcommand{\sectionmark}[1]{}  % prevent sections from overwriting leftmark

% --- Hyperref setup ---
\hypersetup{
  colorlinks=true,
  linkcolor=blue,
  urlcolor=blue,
  citecolor=blue,
}

% --- Caption format ---
\DeclareCaptionFormat{myformat}{\small\bfseries #1#2#3}
\captionsetup{format=myformat}

\begin{document}

\title{AI 系统设计面试题总结}
\date{}
\maketitle
\markboth{TypeScript 版 Agent 知识}{ TypeScript 版 Agent 知识}

\section*{面试官真正想考什么}


AI 系统设计题一般不会满足于“我用 LangChain 搭一个 RAG”。面试官更想看你是否有生产级架构意识。


{\footnotesize
\begin{tabular}{|p{\dimexpr 0.294\columnwidth-2\tabcolsep\relax}|p{\dimexpr 0.392\columnwidth-2\tabcolsep\relax}|p{\dimexpr 0.314\columnwidth-2\tabcolsep\relax}|}
\hline
\textbf{考察方向} & \textbf{面试官想确认什么} & \textbf{常见扣分点} \\
\hline
整体架构 & 你能否把 AI 应用拆成清晰分层 & 上来就讲框架，不讲链路和边界 \\
\hline
模型网关 & 你是否知道模型调用需要统一治理 & 业务代码直接耦合供应商 API \\
\hline
Prompt/Context & 你是否知道提示词和上下文要版本化、可回放 & Prompt 写死在代码里 \\
\hline
RAG/Memory/Tool & 你是否能区分知识、记忆和真实业务动作 & 把所有上下文混在一起塞给模型 \\
\hline
可观测与评测 & 你是否能证明系统质量变化 & 只靠人工试几条问题 \\
\hline
安全合规 & 你是否知道模型不能绕过业务权限 & 只靠 Prompt 防越权和注入 \\
\hline
\end{tabular}
}



系统设计题最怕空泛。好的回答要能沿着一次请求说清楚：用户请求进来后，经过哪些模块，每个模块解决什么问题，出了问题怎么定位，质量下降怎么回滚。



\section*{生产级 AI 应用架构}


参考文章：《AI 应用系统设计：从 Prompt Demo 到生产级架构》



这一组题是 AI 系统设计的核心。你要能把 AI 应用拆成多个工程模块，而不是只说“前端发请求，后端调模型”。



建议掌握这些关键点：


\begin{itemize}[itemsep=1pt, topsep=2pt]
  \item Prompt Demo 证明的是模型能回答，生产系统要证明的是系统能长期、稳定、可控地回答。
  \item 入口层负责鉴权、租户、限流、参数校验和请求分类。
  \item 编排层负责判断任务类型，是普通问答、RAG、Agent、多工具任务，还是异步批处理。
  \item Prompt/Context 层负责模板版本、变量校验、历史消息、检索证据、用户画像和工具说明。
  \item RAG 管共享知识，Memory 管个性化长期事实，Tool 管真实业务动作，三者要分开治理。
  \item 模型网关负责供应商适配、路由、fallback、限流、熔断、Token 预算、成本归因和观测。
  \item 评测观测层负责 Trace、日志、指标、Golden Set、LLM-as-Judge、灰度和回放。
\end{itemize}


高频面试题：


\begin{itemize}[itemsep=1pt, topsep=2pt]
  \item Prompt Demo 到生产系统最大的差距是什么？
  \item 怎么设计一个生产级 AI 应用的整体架构？
  \item 一次 AI 请求从入口到模型返回，完整链路应该怎么讲？
  \item 入口层、编排层、Prompt/Context、RAG/Memory/Tool、模型网关、评测观测分别承担什么职责？
  \item 同步、流式、异步三种模式怎么选？
  \item 为什么需要模型网关？
  \item Prompt 为什么要做版本管理？
  \item RAG 和 Memory 有什么区别？为什么不能混在一起治理？
  \item Tool Calling 的安全边界在哪里？
  \item AI 应用可观测要看哪些指标？
  \item LLM-as-Judge 能不能替代人工评测？
\end{itemize}


回答“怎么设计生产级 AI 应用”时，可以用一个通用模板：先说明业务目标和约束，再讲分层架构，然后讲一次请求链路，接着讲稳定性、安全、成本、观测和评测，最后讲灰度和回滚。这样比直接报一堆技术名词更有说服力。



\section*{稳定性、成本与安全治理}


参考文章：《AI 应用系统设计：从 Prompt Demo 到生产级架构》、《大模型 API 调用工程实践：流式输出、重试、限流与结构化返回》



这一组题考的是生产意识。大模型调用慢、贵、不稳定，输出还不可完全控。没有治理能力，AI 应用很容易在上线后变成成本黑洞和事故来源。



建议掌握这些关键点：


\begin{itemize}[itemsep=1pt, topsep=2pt]
  \item 超时要分层设置：入口超时、模型调用超时、工具调用超时、异步任务超时。
  \item 重试只适合网络瞬断、部分 5xx、供应商过载等可恢复错误；参数错误、权限错误、安全拒答不能盲目重试。
  \item 限流要同时看请求数、Token 数、并发数、租户预算和模型供应商配额。
  \item fallback 要谨慎。模型降级可能影响质量、格式、工具调用能力和安全策略，不是所有任务都能自动降级。
  \item Token 成本要归因到租户、用户、功能、模型、Prompt 版本和业务场景。
  \item Tool Calling 安全必须由后端强制执行，不能相信模型自己判断权限。
\end{itemize}


高频面试题：


\begin{itemize}[itemsep=1pt, topsep=2pt]
  \item AI 应用如何做超时、重试、限流、熔断和降级？
  \item 为什么大模型调用限流要同时看 RPM、TPM、并发数和租户预算？
  \item 如何设计模型 fallback 策略？什么时候不能自动降级？
  \item Token 成本怎么归因到租户、用户、功能和 Prompt 版本？
  \item 高风险工具调用为什么要做二次确认？
  \item PII 脱敏、权限过滤、审计日志应该放在哪些环节？
  \item Prompt 注入攻击在系统设计层面怎么防？
  \item 出现模型输出事故后，如何通过 Trace 回放定位问题？
\end{itemize}


回答安全题时，一定要强调：Prompt 只能辅助，不能替代代码层面的权限校验。模型可以建议调用工具，但后端必须校验用户身份、资源归属、参数范围、操作风险和幂等状态。



\section*{评测与持续迭代}


参考文章：《AI 应用评测体系：从 Golden Set 构建到线上灰度闭环》



传统系统上线前可以跑单元测试、集成测试、压测；AI 应用还要评测答案质量、检索质量、工具轨迹和结构化输出稳定性。没有评测闭环，就很难知道一次 Prompt 调整、模型切换、检索参数变化到底是提升还是退步。



建议掌握这些关键点：


\begin{itemize}[itemsep=1pt, topsep=2pt]
  \item Golden Set 是发布前质量回归的基础，应该覆盖正常路径、边缘场景、对抗样本和高权重失败。
  \item 离线评测适合发布前阻断明显退步，Trace 回放适合复现真实线上路径，线上灰度适合验证真实用户分布。
  \item RAG 要分检索和生成评测，Agent 要看任务完成率、工具选择、参数准确率和轨迹质量。
  \item LLM-as-Judge 可以提高效率，但要用人工抽样、规则校验和参考答案校准。
  \item 评测结果要和 Prompt 版本、模型版本、检索配置、代码版本绑定，便于回滚和定位。
\end{itemize}


高频面试题：


\begin{itemize}[itemsep=1pt, topsep=2pt]
  \item 为什么没有评测集就很难放心上线？
  \item Golden Set 如何覆盖正常路径、边缘场景、对抗样本和高权重失败？
  \item 离线评测、Trace 回放、线上灰度分别放在发布流程的哪个阶段？
  \item RAG、Agent、结构化输出的评测指标为什么不能混用一套？
  \item LLM-as-Judge 有哪些偏差？生产中怎么校准？
  \item CI 自动评测怎么控制成本和耗时？
  \item 线上质量下降时，如何判断是模型、Prompt、检索、工具还是数据分布变化导致？
\end{itemize}


面试里可以把评测讲成一条流水线：开发阶段跑小规模核心 Golden Set，合并或发布前跑完整评测，灰度阶段做线上抽样，事故后用 Trace 回放复现，失败样本再回流到评测集。



\section*{实时语音 Agent}


参考文章：《AI 语音技术详解：从 ASR、TTS 到实时语音 Agent 的工程化落地》



实时语音 Agent 是很典型的 AI 系统设计题，因为它同时考多模态链路、低延迟、状态机、打断处理和端云选型。



建议掌握这些关键点：


\begin{itemize}[itemsep=1pt, topsep=2pt]
  \item 语音 Agent 不是 ASR + LLM + TTS 的简单拼接，而是一套实时音频流系统。
  \item 完整链路包括音频采集、VAD、ASR、LLM、工具调用、TTS、流式播放和打断处理。
  \item 端到端延迟来自多个环节：音频帧提交、VAD 判断、ASR 转写、LLM 首字、TTS 首包、网络和播放缓冲。
  \item 打断处理要取消播放、取消生成、处理已播放内容和未播放内容，并更新对话状态。
  \item 云端 API 上线快，本地模型可控但工程成本高，端云混合更适合兼顾体验和成本。
\end{itemize}


高频面试题：


\begin{itemize}[itemsep=1pt, topsep=2pt]
  \item 如何设计一个实时语音 Agent？
  \item ASR、LLM、TTS、VAD 在语音系统中分别负责什么？
  \item 实时语音 Agent 的端到端延迟主要来自哪里？
  \item 用户打断时，系统应该如何取消播放、取消生成和更新上下文？
  \item listening、thinking、speaking、interrupted 这些状态如何管理？
  \item 云端 API、本地模型、端云混合怎么选？
  \item Speech-to-Speech API 适合什么场景？有哪些取舍？
  \item 语音 Agent 的可观测指标应该包括哪些？
\end{itemize}


回答实时语音题时，可以先拆链路，再讲低延迟优化，接着讲状态机和打断，最后讲可观测和选型。不要只停留在“调用语音识别和语音合成接口”。



\section*{系统设计答题模板}


遇到开放式 AI 系统设计题，可以按下面顺序回答：


\begin{enumerate}[itemsep=1pt, topsep=2pt]
  \item \textbf{明确场景和约束}：用户规模、响应时延、数据来源、权限要求、成本预算、质量目标。
  \item \textbf{拆分核心链路}：入口、编排、上下文、RAG、Memory、Tool、模型网关、输出校验、观测评测。
  \item \textbf{讲关键数据流}：一次请求如何鉴权、检索、组装 Prompt、调用模型、处理流式输出、记录 Trace。
  \item \textbf{补治理能力}：限流、熔断、重试、幂等、fallback、成本归因、权限控制、审计日志。
  \item \textbf{讲评测闭环}：Golden Set、离线评测、Trace 回放、线上灰度、失败样本回流。
  \item \textbf{说明取舍边界}：哪些场景同步，哪些场景流式，哪些场景异步；哪些任务允许降级，哪些必须人工确认。
\end{enumerate}


这套模板能覆盖大多数 AI 应用系统设计题，包括智能客服、企业知识库、代码助手、数据分析 Agent、语音 Agent。



\section*{常见扣分点}

\begin{itemize}[itemsep=1pt, topsep=2pt]
  \item 上来就讲框架名，不讲业务约束和系统边界。
  \item 只讲 Prompt 和模型，不讲 RAG、Memory、Tool 的治理差异。
  \item 没有模型网关意识，业务代码直接调用供应商 API。
  \item 不记录 Prompt 版本、模型版本、检索结果、工具轨迹，导致事故无法回放。
  \item 把 LLM-as-Judge 当成万能评测，不做人工校准和规则校验。
  \item 只靠 Prompt 做安全防护，忽略权限、脱敏、审计和二次确认。
  \item 没有灰度、回滚和失败样本回流机制。
\end{itemize}


\section*{复习建议}


AI 系统设计面试要按“系统链路”来回答，不要从某个框架或工具名开始。更稳的表达方式是先讲 Demo 和生产差距，再讲分层架构、核心链路、治理能力和评测闭环。



如果面试官继续追问，再展开模型网关、Prompt 版本、RAG 和 Memory 隔离、Tool Calling 安全、Trace 回放、灰度评测这些关键点。



最后记住一句话：\textbf{AI 系统设计不是让模型回答一次，而是让系统长期、稳定、可控地回答。} 能把这句话展开成架构、链路、治理和评测，基本就能答到面试官想听的层次。


\end{document}
